\section{The major term of the determinant estimate}\label{sec:major-arcs}\label{maj:section}

We retain the notation of Theorem~\ref{thm:type-II}, in particular
$X=H_mH_n\asymp x$ and $W=\exp(L^{0.24})$.  The exponent $A_0$
defining $\mathfrak M$ is fixed throughout this section.  Constants may
depend on $K$ and on its fixed prime bands; the powers of $L$ used before
choosing $K$ will not depend on $K$.

\begin{proposition}[The determinant major term]\label{prop:major-term}
For every fixed $D>0$, the unrestricted major sum
\begin{equation}\label{maj:sum}
 \begin{split}
 \mathcal Q_Y^{\mathrm{maj}}={}&
 \left(\prod_{i=1}^K V_i^{-2}\right)
 \sum_{a,b,m,n,r,s}
 \eta(a/Y)\eta(b/Y)\alpha_m\overline{\alpha_r}
       \beta_n\overline{\beta_s}\,
 \mathcal H_{\mathfrak M}(bmn-ars;a,b)
 \end{split}
\end{equation}
satisfies
\begin{equation}\label{maj:bound}
 |\mathcal Q_Y^{\mathrm{maj}}|\ll_D XYL^{-D}.
\end{equation}
Here $a,b$ run independently over products of one prime from each
$\mathcal P_i$, with no disjointness condition, and
$\mathcal H_{\mathfrak M}$ denotes the kernel in
\eqref{eq:det-major}.  The estimate is uniform in the coefficient
intervals, in $|v|\le L^C$, and in the admissible scale $Y$.
\end{proposition}

The major-arc decomposition will separate three factors: the centered
prime-minus-rough coefficient $\alpha$, the arbitrary rough coefficient
$\beta$, and the product of small prime labels. Uniform cancellation
of the first factor alone does not control an integral over a frequency
range of length comparable to $X$. We isolate one small prime label.
Where its polynomial is small, a mean-square bound for the product of
the two long polynomials suffices; where it is large, the frequencies
are sparse enough to control the
energy of the $\beta$ polynomial. The next two lemmas supply the
uniform cancellation and the sparse-frequency bound, respectively.
All characters below are extended by zero away from the units.

\begin{lemma}[The long coefficient]\label{maj:long-coefficient}
Fix $A_0,B,A>0$.  For every character $\chi$ of modulus at most
$L^{A_0}$, put
\begin{equation}\label{maj:M-definition}
 M_\alpha(t)=\frac1{H_m}\sum_m\alpha_m\chi(m)m^{it}.
\end{equation}
Then, uniformly for $|t|\le 2XL^B$,
\begin{equation}\label{maj:M-uniform}
 |M_\alpha(t)|\ll_A L^{-A}.
\end{equation}
The assertion also holds with conjugated coefficients and characters.
\end{lemma}

\begin{proof}
Write $u=t+v$ and let $I\subset[H_m,2H_m]$ be the coefficient interval.
We give the approximation of rough numbers explicitly.  Set
$\mathcal P(W)=\{p:p\le W\}$ and
$S_W=\sum_{p\le W}p^{-1}=0.24\log L+O(1)$.
For an integer $r\asymp c\log L$, define
\[
 T_r(m)=\sum_{\substack{d\mid m\,,\ d\ \mathrm{squarefree}\\
                       p\mid d\Rightarrow p\le W\,,\ \#\{p:p\mid d\}\le r}}
             \mu(d).
\]
The two consecutive Bonferroni truncations bracket
$\1_{P^-(m)>W}$.  Their difference is supported on products of $r+1$
distinct primes.  Consequently, on every subinterval $I'\subset I$,
\begin{equation}\label{maj:rough-truncation}
 \sum_{m\in I'}|\1_{P^-(m)>W}-T_r(m)|
 \ll H_m\frac{S_W^{r+1}}{(r+1)!}+W^{r+1},
\end{equation}
where changing $r$ by one, if necessary, is harmless.  This follows
by summing $\1_{d\mid m}$ over the first omitted degree, using
$\#\{m\in I':d\mid m\}=|I'|/d+O(1)$.  Since
$r!\ge(r/e)^r$, choosing $c$ sufficiently large makes the first
term $O(H_mL^{-D'})$ for any prescribed fixed $D'$.  The second is
$x^{o(1)}$, since $r\log W=O(L^{0.24}\log L)$.  Moreover, the
absolute reciprocal coefficient mass of every truncation satisfies
\begin{equation}\label{maj:rough-reciprocal}
 \sum_{d\text{ used}}\frac1d
 \le\prod_{p\le W}(1+p^{-1})\ll L^{0.24}.
\end{equation}
In particular, increasing the depth to improve the error does not
increase the exponent in this bound.

First consider large $|u|$.  For each surviving divisor $d$ with
$(d,k)=1$, where $k$ is the character modulus, write $m=dz$ and
split $z$ into residues modulo $k$.  Its scale is
$N'=H_m/d=x^{\Omega(1)}$.  For
$1\le |u|\le T_*=\exp(L/(\log L)^2)$, summation against an integral
on a progression gives
\begin{equation}\label{maj:phase-elementary}
 \sum_{\substack{z\in J\subset[N',2N']\\z\equiv a\pmod k}}z^{iu}
  =\frac1k\int_J y^{iu}\,\dd y+O(1+|u|),
 \qquad
 \left|\int_Jy^{iu}\,\dd y\right|\ll\frac{N'}{|u|}.
\end{equation}
Indeed the total variation of $y^{iu}$ on this dyad is $O(|u|)$;
the integral formula follows by evaluating $y^{1+iu}/(1+iu)$ at
the endpoints.  The error, after summing the at most $k$ residues
and all $d\le W^r$, is $x^{-c_\delta}$ after division by $H_m$,
uniformly up to $T_*$.  At $T_*\le |u|<x^2$, apply
Lemma~\ref{lem:log-phase}: its lower scale condition holds because
$N'\ge x^{\delta/2}$ for sufficiently large $x$, and its lower
frequency condition is weaker than $|u|\ge T_*$.  Together with
\eqref{maj:rough-reciprocal}, these bounds and partial summation
for $1/\log m$ imply
\begin{equation}\label{maj:rough-high}
 \begin{split}
 \left|\frac1{H_mV(W)}
       \sum_{\substack{m\in I\\P^-(m)>W}}
       \frac{\chi(m)m^{iu}}{\log m}\right|
 \ll{}& L^{-D'}+
 \begin{cases}
 L^{C_0}(|u|^{-1}+x^{-c_\delta}),&1\le |u|\le T_*,\\
 L^{C_0}\exp(-L/(\log L)^{C_3}),&T_*\le |u|<x^2.
 \end{cases}
 \end{split}
\end{equation}
Here $C_0$ is fixed after $A_0,\delta$; it need not grow with the
chosen Bonferroni depth.  The factor $1/V(W)\ll L^{0.24}$ and the
bounded variation of $1/\log m$ have been included in $C_0$.

For the prime part, choose fixed $0<\tau<\delta$ and
$1-\delta<\eta<1$.  The scale conditions imply
$x^\tau/2\le H_m\le x^\eta$ for large $x$.
Lemma~\ref{lem:long-prime}, followed by partial summation to replace
$p^{-1}$ by $H_m^{-1}$, gives an arbitrary negative power of $L$
when $L^{B_0'}\le |u|\le x^2$.  Choose $B_2>C+2$ sufficiently
large after $A,A_0,B$ and the exponent $C_0$.  Then
$L^{B_2}\le |t|\le2XL^B$ implies
$|u|\asymp|t|$, $|u|\ge L^{B_0'}$, and $|u|<x^2$.
The prime estimate and \eqref{maj:rough-high} prove
\eqref{maj:M-uniform} on this range.

It remains to treat $|t|\le L^{B_2}$, hence $|u|\le2L^{B_2}$.
Choose the counting precision $D'$ after $B_2$.
Lemma~\ref{lem:rough-counts} gives, on every subinterval,
\begin{equation}\label{maj:rough-low-count}
 \sum_{m\in I'}\chi(m)\1_{P^-(m)>W}
 =\1_{\chi\ \mathrm{principal}}V(W)|I'|
       +O(H_mL^{-D'}).
\end{equation}
The principal character equals one on rough numbers because all
primes dividing $k$ are at most $W$; for nonprincipal characters
the count has zero main term by complete-period cancellation.
For primes, Lemma~\ref{lem:prime-estimates} and character
orthogonality give, to any prescribed precision,
\[
 \sum_{p\in I'}\chi(p)
 =\1_{\chi\ \mathrm{principal}}\int_{I'}\frac{\dd y}{\log y}
       +O(H_mL^{-D'}).
\]
Partial summation introduces at most a fixed power of $L$ on this
low-frequency range.  Dividing the rough main term by
$V(W)\log y$ produces exactly $\dd y/\log y$; its twisted main
term therefore cancels the prime main term.  Increasing $D'$ gives
\eqref{maj:M-uniform}.  Complex conjugation changes only the signs
of the real frequencies and the character, so the proof is unchanged.
\end{proof}

\begin{lemma}[A small prime factor and sparse frequencies]
\label{maj:sparse-frequency}
Let $P$ satisfy
$cL^{0.1}\le\log P\le C_1L^{0.2}$, and let $\mathcal P$ be any
set of primes in $[P,2P]$.  For any real $\xi$ and any character
$\chi$, put
\[
 P_s(t)=P^{-1}\sum_{p\in\mathcal P}\chi(p)p^{i(t+\xi)}.
\]
Let
$N_\beta(t)=H_n^{-1}\sum_n\beta_n\chi_n(n)n^{it}$, with
$|\beta_n|\le L^C$ on $[H_n,2H_n]$ and
$x^\delta\le H_n\ll x$.
For every fixed $B,A_s>0$, the unit intervals meeting
\[
 \{t:|t|\le 2XL^B,\ |P_s(t)|>L^{-A_s}\}
\]
form a family $\mathcal I$ satisfying
\begin{equation}\label{maj:sparse-count}
 |\mathcal I|\le\exp(O(L^{0.91})),
 \qquad
 \sum_{I\in\mathcal I}\sup_{t\in I}|N_\beta(t)|^2
       \ll L^{2C+1}.
\end{equation}
Both assertions are uniform in $\xi$ and in the prime set.
\end{lemma}

\begin{proof}
Put $Z=4XL^B$ and $j=\lfloor\log Z/\log(2P)\rfloor$.
Write $P_s(t)^j=\sum_{n\le(2P)^j}c_nn^{it}$.
For each product $n$, unique factorization bounds the number of
ordered prime tuples producing it by $j!$.  Hence
\begin{equation}\label{maj:factorial-norm}
 \sum_n|c_n|^2
 \le j!P^{-2j}|\mathcal P|^j
 \le j!(C_2/P)^j.
\end{equation}
This is a coefficient estimate for the actual growing degree $j$;
no fixed-divisor-moment constant is used for it.

Choose an exceeding point from each member of $\mathcal I$, and
split these points into three classes according to the integer
left endpoint modulo three.  Each class is separated by at least
one.  The separated mean-square estimate of
Lemma~\ref{lem:moment-bounds}, whose constant is independent of
the coefficients, gives
\begin{equation}\label{maj:exception-count-calculation}
 |\mathcal I|
 \ll ZL^{O(1)}L^{2A_sj}j!(C_2/P)^j.
\end{equation}
Since $\log Z-j\log P\ll\log P+j$ and
$j\ll L^{0.9}$, the logarithm of its right-hand side is at most
\[
 O(\log P+j\log(j+1)+A_sj\log L+\log L)
       =O(L^{0.9}\log L+L^{0.2})=O(L^{0.91}).
\]
The factor $p^{i\xi}$ affects no coefficient absolute value, so
this estimate is uniform for all real $\xi$.

For the second assertion choose a maximizing point $t_I$ on the
closure of each unit interval, and again use three separated
classes.  In one class consider the evaluation matrix on the full
integer interval $[H_n,2H_n]$.  Its Gram entries are
\[
 G_{I,J}=\sum_{H_n\le n\le2H_n}n^{i(t_I-t_J)}.
\]
Set $T_*=\exp(L/(\log L)^2)$.  The diagonal is $O(H_n)$.
For $1\le|\Delta|\le T_*$,
\eqref{maj:phase-elementary} with modulus one gives
\begin{equation}\label{maj:Gram-near}
 H_n^{-1}|G_{I,J}|\ll |\Delta|^{-1}+x^{-c_\delta}.
\end{equation}
For $T_*<|\Delta|\le 4XL^B+2<4x^3$,
Lemma~\ref{lem:log-phase} gives
\begin{equation}\label{maj:Gram-far}
 H_n^{-1}|G_{I,J}|
 \ll\exp(-L/(\log L)^{C_3}).
\end{equation}
All its scale hypotheses hold because $H_n\ge x^\delta$ and
$H_n\ll x$.  Separation implies that the sum of $1/|\Delta|$
in any row is $O(\log(2Z))=O(L)$.
The other row sums tend to zero, since
\[
 |\mathcal I|x^{-c_\delta}=o(1),\qquad
 |\mathcal I|\exp(-L/(\log L)^{C_3})=o(1).
\]
Thus the absolute Gram row sums are $O(H_nL)$, and the same is
true of its operator norm.  The coefficient vector of $N_\beta$
has squared norm $O(L^{2C}/H_n)$.  Applying the Gram bound and
adding the three classes proves \eqref{maj:sparse-count}.
\end{proof}

\begin{proof}[Proof of Proposition~\ref{prop:major-term}]
For large $x$ the arcs are disjoint: distinct reduced rationals of
denominator at most $L^{A_0}$ are at distance at least $L^{-2A_0}$,
whereas their radii are $2L^{A_0}/Y$ and $Y\ge\exp(cL^{0.1})$.
On an arc write $\theta=h/k+\lambda/Y$, with
$|\lambda|\le2L^{A_0}$ and $\dd\theta=\dd\lambda/Y$.
Every variable in \eqref{maj:sum} is a unit modulo $k$: the group
primes and all prime factors of the rough variables exceed $k$.
The function
\[
 (a,b,m,n,r,s)\longmapsto
 \e\big(h(bmn-ars-b+a)/k\big)
\]
on $((\Z/k\Z)^\times)^6$ has an exact Fourier expansion in
products of six multiplicative characters.  Each Fourier
coefficient has absolute value at most one, so their total absolute
mass is at most $\varphi(k)^6\le k^6$.  There are $O(L^{2A_0})$
arcs.  It is therefore enough to estimate each separated character
term, uniformly in $\lambda$, with arbitrarily large logarithmic
saving.

Set $w=bmn/(YX)$ and $w'=ars/(YX)$.  Both lie in a fixed compact
subinterval of $(0,\infty)$.  Apart from the factors
$\e(-\lambda b/Y)\e(\lambda a/Y)$, the remaining common kernel is
\[
 \psi_\lambda(X(w-w')),\qquad
 \psi_\lambda(y)=\psi(y)\e(\lambda y).
\]
Here is the Mellin separation with its normalization.  In coordinates
$z=\log w'$ and $\zeta=X\log(w/w')$, it is
\[
 \psi_\lambda\big(Xe^z(e^{\zeta/X}-1)\big).
\]
Insert fixed smooth cutoffs equal to one on the possible values of
$w,w'$.  The resulting function is supported on a fixed compact
set of $(\zeta,z)$: the support of $\psi$ forces $|\zeta|\ll1$.
Its derivatives of order $r+s$ are
$O_{r,s}(L^{B_0(r+s)})$ for some fixed $B_0$ chosen after $A_0$.
Fourier inversion, with inverse kernel $e^{i(v\zeta+s'z)}$, and
the change of variable $t=Xv$ consequently give the exact identity
\begin{equation}\label{maj:Mellin-separation}
 \psi_\lambda(X(w-w'))
 =\frac1X\iint_{\R^2}
      F_\lambda(t/X,s')w^{it}(w')^{i(s'-t)}\,\dd t\,\dd s'
\end{equation}
on the coefficient support.  Repeated integration by parts yields,
for every fixed integer $j\ge1$,
\begin{equation}\label{maj:transform-decay}
 |F_\lambda(v,s')|
 \ll_j(1+|v|/L^{B_0})^{-j}(1+|s'|/L^{B_0})^{-j}.
\end{equation}

Define the normalized endpoint polynomials
\begin{equation}\label{maj:polynomials}
 \begin{split}
 B_\lambda(t)&=\frac1Y\sum_b\eta(b/Y)\e(-\lambda b/Y)
                 \left(\prod_iV_i^{-1}\right)\chi_b(b)b^{it},\\
 N_\beta(t)&=\frac1{H_n}\sum_n\beta_n\chi_n(n)n^{it},
 \qquad Q(t)=B_\lambda(t)M_\alpha(t)N_\beta(t).
 \end{split}
\end{equation}
The other endpoint gives a polynomial $Q'$ of the same type with
conjugations and sign changes.  Each unnormalized endpoint sum is
$YX$ times its polynomial.  Combining these two factors with the
$X^{-1}$ in \eqref{maj:Mellin-separation} and the arc measure $Y^{-1}$
gives
\begin{equation}\label{maj:XY-normalization}
 (YX)^2\cdot X^{-1}\cdot Y^{-1}=XY.
\end{equation}
The additional factor $(YX)^{-is'}$ has absolute value one.
Thus, after division by $XY$, the separated expression is an
integral of $F_\lambda(t/X,s')Q(t)Q'(s'-t)$, followed by the
$\lambda$ integral.

The bound $|B_\lambda(t)|\ll1$ follows directly from
$b\asymp Y$ and
\[
 \sum_b\frac1b\prod_iV_i^{-1}
 \le\prod_i\left(V_i^{-1}\sum_{p\in\mathcal P_i}p^{-1}\right)=1.
\]
The coefficient of index $u=mn$ in $M_\alpha N_\beta$ is
$O(L^C\tau(u)/X)$ and is supported on $u\asymp X$.
Lemma~\ref{lem:moment-bounds} therefore gives a squared coefficient
norm $O(L^{C_4}/X)$ and, on every real interval $I$ of length
$Z\ge1$,
\begin{equation}\label{maj:product-mean}
 \int_I|M_\alpha(t)N_\beta(t)|^2\,\dd t
       \ll L^{C_4}(1+Z/X).
\end{equation}
The fixed exponent $C_4$ depends only on the original coefficient
bounds.  The estimate is uniform in translates of $I$, by twisting
the coefficients, and also bounds the mean square of $Q$ and $Q'$.

Fix $B_1>B_0+1$.  Equations \eqref{maj:transform-decay} and
\eqref{maj:product-mean} permit us, to arbitrary logarithmic
precision, to restrict to
\begin{equation}\label{maj:retained-rectangle}
 |s'|\le L^{B_1},\qquad |t|\le XL^{B_1}.
\end{equation}
For completeness, on a dyadic $t$ interval of length $Z$,
Cauchy's inequality bounds the integral of
$|Q(t)Q'(s'-t)|$ by $L^{C_4}(1+Z/X)$, uniformly in $s'$.
For $Z\ge XL^{B_1}$ multiply this by
$(1+Z/(XL^{B_0}))^{-j}$ and sum the geometric tail.
For the $s'$ tail, the full weighted $t$ integral is
$O(L^{C_4+B_0+1})$, uniformly in $s'$, and
\[
 \int_{|s'|>L^{B_1}}(1+|s'|/L^{B_0})^{-j}\,\dd s'
 \ll_j L^{B_0-(B_1-B_0)(j-1)}.
\]
Taking $j$ large makes both errors as small as required, including
the character and arc costs.  This argument also covers translated
frequencies $s'-t$ without a pointwise tail assumption.

On \eqref{maj:retained-rectangle}, Cauchy's inequality in $t$
reduces the claim to proving, for every prescribed $A'>0$,
\begin{equation}\label{maj:energy}
 \int_{|t|\le2XL^{B_1}}
       |B_\lambda(t)M_\alpha(t)N_\beta(t)|^2\,\dd t
       \ll_{A'}L^{-A'}.
\end{equation}
Indeed the retained $s'$ integral has length $2L^{B_1}$ and all
remaining character and $\lambda$ costs are fixed powers of $L$.

Write $b=pb'$ with $p\in\mathcal P_1$ and split this prime band
into $O(L)$ ordinary dyads $[P,2P)$.  The cutoff on $b$ permits
restriction to $Y/(2P)\le b'\le4Y/P$.  Mellin inversion of
$f_\lambda(y)=\eta(y)\e(-\lambda y)$ gives
\[
 f_\lambda(y)=\int_\R\widehat f_\lambda(\xi)y^{i\xi}\,\dd\xi,
 \qquad \int_\R|\widehat f_\lambda(\xi)|\,\dd\xi
                \ll L^{C_5},
\]
where $C_5$ is fixed after $A_0$.  This last bound follows, for
example, by twice integrating by parts outside $|\xi|\le1$;
the first two derivatives of $f_\lambda(e^z)$ have fixed compact
support and size $O((1+|\lambda|)^2)$.
The remaining factor
\[
 Q_{b'}(u)=\frac P Y\sum_{Y/(2P)\le b'\le4Y/P}
      \left(\prod_{i=2}^KV_i^{-1}\right)\chi_b(b')(b')^{iu}
\]
has absolute value $O(1)$ by its reciprocal coefficient mass.
Precisely, the dyad contribution to $B_\lambda(t)$ is
\[
 V_1^{-1}\int_\R\widehat f_\lambda(\xi)Y^{-i\xi}
     \left(P^{-1}\sum_{p\in\mathcal P_1\cap[P,2P)}
                         \chi_b(p)p^{i(t+\xi)}\right)
                  Q_{b'}(t+\xi)\,\dd\xi.
\]
Minkowski's integral inequality shows that it suffices to prove
\begin{equation}\label{maj:isolated-energy}
 \int_{|t|\le2XL^{B_1}}
       |P_s(t)M_\alpha(t)N_\beta(t)|^2\,\dd t\ll L^{-A''}
\end{equation}
with arbitrary fixed $A''$, uniformly in every real shift $\xi$.
The $O(L)$ dyads and the Mellin $L^1$ norm are then paid by
increasing $A''$.  In particular there is no unaccounted Fourier
tail in this factor separation.

On $|P_s(t)|\le L^{-A_s}$, Equation \eqref{maj:product-mean}
bounds \eqref{maj:isolated-energy} by
$O(L^{-2A_s+C_4+B_1})$.  Choose $A_s$ large after $A''$.
On the remaining set use the unit intervals of
Lemma~\ref{maj:sparse-frequency}.  We have $|P_s(t)|\ll1$, and
Lemma~\ref{maj:long-coefficient} gives
$|M_\alpha(t)|\ll L^{-A}$ everywhere in the retained range,
including its low frequencies.  Therefore its contribution is
\[
 \ll L^{-2A}
       \sum_{I\in\mathcal I}\sup_{t\in I}|N_\beta(t)|^2
 \ll L^{-2A+2C+1}.
\]
Taking $A$ sufficiently large proves
\eqref{maj:isolated-energy}, then \eqref{maj:energy}, and finally
\eqref{maj:bound}.
\end{proof}

\begin{proof}[Completion of Theorem~\ref{thm:type-II}]
Fix the requested precision $D_*$ and choose the expanded-square
precision $D_1$ sufficiently large in terms of $D_*$ and the
original coefficient bound $C$.  These choices precede $K$.
The squarefree-label and coincident-label errors from the Cauchy
reduction are smaller than every fixed logarithmic power.
Proposition~\ref{prop:minor-moment}, including its restoration of
good states and harmonic pads, replaces the remaining determinant
equality by \eqref{eq:det-major} with error
$O(XYL^{-D_1})$.

To record the parameter order, first choose $E_0$ after $D_1,C$,
then the arc exponent $A_0$, and finally $K$, as in that
proposition.  For a pad dyad, division of its pairing bound by
$d_0$ uses $UV=d_0YX$ and gives
$XYL^{-E_0+O_C(1)}$.  There are $O_K(L)$ pad dyads; thus $E_0$
can absorb this loss with an exponent independent of $K$.
Constants involving the now fixed $K$ and its band exponents only
affect the threshold for $x$.
The analytic accuracies in Proposition~\ref{prop:major-term} are
chosen last, after $A_0$.  That proposition bounds the unrestricted
major term by $O(XYL^{-D_1})$.

Consequently the Cauchy square in each $Y$ dyad is
$O(XYL^{-D_1})$.  Multiplication by its Cauchy factor $O(X/Y)$
and taking square roots gives $O(XL^{-D_1/2})$ for that dyad.
The $O_K(L)$ dyads of $Y$, and all preceding fixed coefficient
losses, are absorbed by the original choice of $D_1$.
This proves \eqref{eq:type-II}, with $K$ depending on the stated
fixed data and not on the particular exponents $a_i$.
\end{proof}
